\chapter{Devices and I/O}

\section{Devices and buses}

Devices (disks, network cards, timers, serial consoles, GPUs) connect to the CPU and memory
through one or more \textbf{buses}. A typical system has a fast memory bus, an I/O bus
(PCI Express), and slower peripheral buses (USB, SATA). In System/161 all devices sit on a
simple \code{LAMEbus} with up to 32 slots.

\begin{definition}[Device driver]
A \textbf{device driver} is the part of the kernel that interacts with a particular device: it
knows the device's registers and protocol, and presents a uniform interface (e.g.\ ``read block
$n$'') to the rest of the kernel.
\end{definition}

\section{Device registers}

The kernel communicates with a device through \textbf{device registers}, typically of three
kinds:
\begin{description}
  \item[Status] read to learn the device's state (busy, done, error).
  \item[Command] written to tell the device what to do.
  \item[Data] used to transfer data to/from the device.
\end{description}

Two ways to access device registers:
\begin{itemize}
  \item \textbf{Port-mapped I/O}: special instructions (\code{in}/\code{out} on x86) address a
    separate I/O address space.
  \item \textbf{Memory-mapped I/O}: registers appear at physical addresses; ordinary loads and
    stores access them. MIPS/System/161 uses this: each LAMEbus slot gets a 64\,KB region of
    physical addresses starting at \code{0x1FE00000}, which the kernel accesses through the
    uncached segment \code{kseg1} (\code{0xA0000000} + physical) so that reads/writes really reach
    the device.
\end{itemize}

\begin{example}[System/161 timer/clock device]
\begin{center}
\begin{tabular}{@{}cccl@{}}
\toprule
Offset & Size & Type & Description \\
\midrule
0 & 4 & status & current time (seconds) \\
4 & 4 & status & current time (nanoseconds) \\
8 & 4 & command & restart-on-expiry flag \\
12 & 4 & status and command & interrupt (reading acknowledges it) \\
16 & 4 & command & countdown time (microseconds) \\
20 & 4 & command & speaker (causes a beep) \\
\bottomrule
\end{tabular}
\end{center}
Writing 1000000 to offset 16 with restart set makes the timer interrupt every second.
\end{example}

\section{Polling vs interrupts}

\begin{example}[System/161 serial console, write one character]
Registers: \emph{character buffer} (data), \emph{write IRQ} (status/command), \emph{read IRQ}.

\textbf{Polling} (busy waiting):
\begin{lstlisting}[style=c]
void serial_putc_polled(char c) {
    write_reg(SERIAL_CHAR, c);            /* start the operation */
    while (read_reg(SERIAL_WIRQ) == 0)    /* spin until device says done */
        ;
    write_reg(SERIAL_WIRQ, 0);            /* acknowledge completion */
}
\end{lstlisting}
\textbf{Interrupt-driven}:
\begin{lstlisting}[style=c]
void serial_putc(char c) {
    P(serial_write_sem);                  /* only one write in progress (init 1) */
    write_reg(SERIAL_CHAR, c);            /* start; return immediately */
}
void serial_interrupt_handler(void) {     /* device raises an interrupt when done */
    write_reg(SERIAL_WIRQ, 0);            /* acknowledge */
    V(serial_write_sem);                  /* allow the next write */
}
\end{lstlisting}
\end{example}

\begin{important}[Trade-off]
Polling wastes CPU while the device works, but has very low latency and no interrupt overhead;
it is fine for very fast devices or when the CPU has nothing else to do. Interrupts let the CPU
run other threads while the device works, at the cost of interrupt handling overhead per event.
Very high-rate devices (fast network cards) switch adaptively between the two (e.g.\ Linux
NAPI).
\end{important}

\section{Moving data: programmed I/O vs DMA}

\begin{description}
  \item[Programmed I/O (PIO)] The CPU itself copies data between memory and the device's data
    registers or on-device buffer, word by word. Simple, but the CPU is busy for the whole
    transfer.
  \item[Direct memory access (DMA)] The device transfers data directly to/from main memory.
    The driver tells the device the physical memory address, the length and the operation;
    the device then masters the bus, copies the data, and raises an interrupt on completion.
    The CPU is free during the transfer (it competes with the device for the memory bus only).
\end{description}

\begin{lstlisting}[style=c]
/* DMA disk read, conceptually */
void disk_read_dma(uint32_t sector, paddr_t buf, size_t nsectors) {
    lock_acquire(disk_lock);
    write_reg(DISK_SECTOR, sector);
    write_reg(DISK_ADDR,   buf);          /* PHYSICAL address of the buffer */
    write_reg(DISK_COUNT,  nsectors);
    write_reg(DISK_CMD,    CMD_READ);     /* start */
    P(disk_done);                         /* block; interrupt handler does V */
    lock_release(disk_lock);
}
\end{lstlisting}
\begin{remark}
DMA uses physical addresses, so the buffer must be contiguous in physical memory (or the device
must support scatter/gather lists) and must not be paged out during the transfer
(\emph{pinned}).
\end{remark}

System/161 disks use PIO through a 512-byte on-device transfer buffer that is memory-mapped.

\section{Hard disk drives}

\begin{center}
\begin{tikzpicture}
  \draw (0,0) circle (2cm);
  \draw (0,0) circle (1.4cm);
  \draw (0,0) circle (0.8cm);
  \fill (0,0) circle (0.08cm);
  \draw[fill=gray!30] (0:1.4) arc (0:30:1.4) -- (30:2) arc (30:0:2) -- cycle;
  \node[font=\small] at (15:2.5) {sector};
  \node[font=\small] at (-90:1.1) {track};
  \draw[thick] (4,-1.8) -- (1.2,0.6);
  \fill (1.2,0.6) circle (0.07);
  \node[font=\small,align=left,anchor=north west] at (4.05,-1.8) {arm with\\read/write head};
  \draw[-{Stealth}] (130:2.3) arc (130:170:2.3);
  \node[font=\small,anchor=east] at (150:2.45) {rotation};
\end{tikzpicture}
\end{center}

A disk has one or more \textbf{platters} rotating at constant speed (5400--15000 RPM). Each
surface is divided into concentric \textbf{tracks}, each track into \textbf{sectors} (512\,B or
4\,KB). Tracks at the same radius on all surfaces form a \textbf{cylinder}. An \textbf{arm}
moves the read/write heads radially. The drive presents a linear array of logical block numbers;
nearby block numbers are physically nearby.

Cost of an I/O request of $k$ sectors:
\[
  T = T_{\text{seek}} + T_{\text{rot}} + T_{\text{transfer}}.
\]
\begin{description}
  \item[Seek time] move the heads to the right track: from 0 (same track) up to a maximum of
    roughly 10--20\,ms; a common model is linear in seek distance, average $\approx$ one third of
    the maximum for random requests.
  \item[Rotational latency] wait for the sector to come under the head: between 0 and one full
    rotation $\omega^{-1}$; average half a rotation.
  \item[Transfer time] $k/(\text{sectors per track}) \times$ rotation time.
\end{description}

\begin{example}
10000 RPM $\Rightarrow$ one rotation $= 60/10000$\,s $= 6$\,ms, average rotational latency
3\,ms. Average seek 5\,ms. 1000 sectors per track; reading 8 sectors (4\,KB): transfer
$= 8/1000 \times 6\,\text{ms} = 0.048$\,ms. Total $\approx 8.05$\,ms, so random 4\,KB reads
achieve about $4\,\text{KB}/8\,\text{ms} \approx 0.5$\,MB/s, while sequential reads approach
$1000 \times 512\,\text{B} / 6\,\text{ms} \approx 85$\,MB/s.
\end{example}

\begin{important}[Moral]
Seek and rotation dominate. Large, sequential transfers are vastly faster than small random ones.
File systems and disk schedulers try to make I/O sequential and to reduce seek distances.
\end{important}

\section{Disk head scheduling}

When several requests are pending, the order matters. Example: head at cylinder 53 (moving
toward higher numbers), queue $98, 183, 37, 122, 14, 124, 65, 67$, cylinders $0..199$.

\begin{description}
  \item[FCFS] Serve in arrival order. Fair but poor: $53\to98\to183\to37\to122\to14\to124\to65
    \to67$, total movement $45+85+146+85+108+110+59+2 = 640$ cylinders.
  \item[SSTF (shortest seek time first)] Always serve the closest request:
    $53\to65\to67\to37\to14\to98\to122\to124\to183$, total $12+2+30+23+84+24+2+59 = 236$.
    Greedy and good, but can \textbf{starve} far requests.
  \item[SCAN (elevator)] Move in one direction serving requests, reverse at the last request
    in that direction (strictly this variant is called LOOK): $53\to65\to67\to98\to122\to124
    \to183\to37\to14$, total $130 + 169 = 299$. No starvation; middle cylinders get better service.
  \item[C-SCAN] Serve only while moving in one direction, then jump back to the start
    (circular). More uniform waiting times.
\end{description}

Modern drives do their own scheduling internally (they know the exact geometry and rotational
position); the OS keeps queues deep enough for the drive to optimize, and may still merge and
sort requests.

\section{Solid state drives (SSDs)}

SSDs use NAND flash: no moving parts, so no seek or rotational latency; random reads are fast
($\sim$tens of $\mu$s).

\begin{itemize}
  \item Flash is organized into \textbf{pages} (e.g.\ 4\,KB, the unit of read/write) grouped into
    \textbf{blocks} (e.g.\ 128 pages, the unit of erase).
  \item A page cannot be overwritten in place: it must be erased first, and erase works only on
    whole blocks and is slow (ms).
  \item Each block survives a limited number of program/erase cycles (\textbf{wear out}).
  \item The \textbf{flash translation layer} (FTL) in the drive maps logical blocks to physical
    pages. Writes go to fresh erased pages (log-structured style); old copies are marked
    invalid; \textbf{garbage collection} copies valid pages out of mostly-invalid blocks and
    erases them; \textbf{wear levelling} spreads erases evenly.
  \item The OS can help with \textbf{TRIM}: tell the drive which logical blocks are no longer in
    use. Disk head scheduling is irrelevant for SSDs; defragmentation is harmful (extra writes).
\end{itemize}

\section{Other persistent storage}
Persistent memory (e.g.\ byte-addressable NVM) blurs the line between memory and storage; RAID
combines multiple disks for performance (striping) and reliability (mirroring, parity).

\section*{Summary}
\begin{itemize}
  \item Drivers talk to devices through status/command/data registers, via port-mapped or
    memory-mapped I/O (System/161: memory-mapped, accessed through uncached kseg1).
  \item Completion is detected by polling or interrupts; data is moved by programmed I/O or DMA.
  \item Disk access time $=$ seek $+$ rotational latency $+$ transfer; seek and rotation dominate.
  \item Disk scheduling: FCFS (fair, slow), SSTF (greedy, starvation), SCAN/elevator and C-SCAN.
  \item SSDs: no mechanical latency, but erase-before-write in blocks, wear; the FTL hides this.
\end{itemize}
